% Capítulo 5 - Conclusiones y trabajo futuro.
% Fuentes: claude/ideas-pendientes.md y la sección 10 del tutorial.
% Extensión objetivo: 3-5 páginas.

\chapter{Conclusiones y trabajo futuro}
\label{ch:conclusiones}

\section{Conclusiones}
\label{sec:conclusiones}

% TODO: redactar. Estructura sugerida:
% 5.1.1 Cumplimiento de objetivos.
%   - Repaso de los O1..O6 del capítulo 1 con marca de cumplido/parcial
%     y referencia al apartado donde se justifica.
% 5.1.2 Aportaciones principales.
%   - Modelo de árbol con herencia en cascada y cierre léxico de
%     plantillas.
%   - Patrón cliente/servidor con frontera REST en una imagen Docker
%     descargable.
%   - Disciplina de calidad: ~95% cobertura, mypy estricto, CI/CD.
%   - Comparativa con esxobjects/yamlinfr como base reutilizable.
%   - Metodología de trabajo con asistencia IA, trazable y
%     reproducible.
% 5.1.3 Limitaciones conocidas (honestas).
%   - El conector real (VMware/Proxmox) no está implementado en esta
%     versión.
%   - La interfaz web no entra en esta versión.
%   - La aplicación específica para docencia (capa 2) está diseñada
%     pero no implementada.
%   - No hay autenticación: el servicio se asume desplegado en red de
%     confianza.

\section{Trabajo futuro dentro del TFG}
\label{sec:futuro-tfg}

% TODO: redactar. Las tres piezas que sí entran en el TFG:
% - Conectores reales contra VMware del aulario (acceso por VPN
%   pendiente del tutor).
% - Interfaz web básica con NiceGUI sobre la API REST: navegador del
%   árbol, edición por diálogos, vista alternativa "desde el
%   hipervisor".
% - Aplicación específica para despliegue de asignaturas (capa 2): una
%   única orden despliega toda la asignatura a partir de una lista de
%   alumnos y una plantilla común.

\section{Trabajo futuro más allá del TFG}
\label{sec:futuro-mas-alla}

% TODO: redactar como lista, con párrafo introductorio que justifique
% por qué quedan fuera (alcance temporal del TFG). Lista:
% - Control de usuarios, autenticación y permisos (RBAC). Diseño
%   completo cerrado, puntos de extensión identificados.
% - Interfaz web más completa: acciones, edición avanzada, dashboards
%   operacionales.
% - Integración con sistemas de gestión de direcciones IP (IPAM).
% - Gestión de snapshots de VMs con políticas heredables.
% - Modelado de datastores como recurso heredable.
% - Conectores adicionales: Hyper-V, KVM, plataformas cloud.
% - Alta disponibilidad del servicio (activo/pasivo, MongoDB en
%   sustitución de TinyDB).
% - Despliegue de la configuración desde repositorios git (GitOps).
% - Logs estructurados a syslog (módulo logging stdlib, ya diseñado en
%   DEC-023 pero pendiente de implementar).
% - Política de seguridad de undeploy con flag --force.
